\section{Discussion}

% ==========================================
% 4. SECTION 4.1 
% ==========================================

\subsection{The Role of Explicit Inductive Bias in Small-Scale Molecular Datasets}

Graph Neural Networks theoretically offer superior representational capacity by learning molecular features directly from graph topology, thereby bypassing the need for manual feature engineering. However, this architectural flexibility demands extensive training data to converge on generalizable chemical representations. In the context of the TgDHFR dataset, which comprises only \NCompounds\ compounds, the GNN architecture failed to leverage its theoretical superiority, plateauing at a predictive performance statistically indistinguishable from classical machine learning approaches. 

This convergence failure highlights the critical role of explicit inductive bias in low-data regimes. The fundamental lack of explicit chemical priors in deep representation learning forces these architectures to construct a continuous latent space entirely from scratch. With fewer than 1,000 training instances, the resulting latent space density is severely restricted, creating a sparse manifold where interpolation between data points becomes highly unstable. To counteract this sparsity, deep learning frameworks structurally require artificial variance injection. Consequently, achieving high predictive accuracy with neural networks on this specific dataset necessitates extensive Gaussian noise-based data augmentation, as demonstrated in the original study \cite{karimova2025}, to artificially inflate latent space density and force convergence.

Classical tree-based ensembles operating on pre-computed Morgan fingerprints bypass the need for end-to-end representation learning by utilizing fixed, domain-specific descriptors. These fingerprints encode established chemical heuristics—specifically, circular topological substructures—directly into the feature space. For small-scale datasets, this explicit pre-computation acts as a natural regularizer, rigidly constraining the hypothesis space and preventing the model from memorizing noise. Our results explicitly demonstrate that exploiting the inductive bias inherent to traditional structural fingerprints entirely eliminates the necessity for computationally burdensome data augmentation in this specific context. 

Nevertheless, this architectural superiority is strictly regime-dependent. While deterministic feature engineering provides a parsimonious and highly reliable strategy for the current low-data regime ($N = \NCompounds$), this advantage is expected to invert as the dataset scales. In high-data environments ($N \gg 1000$), the rigid constraints of pre-computed fingerprints typically become an informational bottleneck. Under those conditions, the flexible representational capacity of GNNs allows them to learn task-specific embeddings that systematically outperform classical descriptors.

% ==========================================
% 4. SECTION 4.2
% ==========================================

\subsection{Dimensionality, Conformational Overhead, and Model Efficiency}

Computational efficiency determines the practical viability of molecular representations in high-throughput virtual screening. Generating 3D geometric descriptors requires extensive conformational sampling and energy minimization (e.g., ETKDGv3 embedding and MMFF94s optimization). This conformational overhead scales non-linearly, creating severe processing bottlenecks when screening libraries expand into the millions of compounds.

Empirical results demonstrate that this computational expenditure provides no predictive advantage for TgDHFR inhibition. The composite 2D/3D/FP representation yielded an $R^2$ of \RdkitTwoDThreeDFpRTwoMean, which remains statistically indistinguishable from the purely topological Morgan fingerprint model ($R^2 = \MorganRTwoMean$; Nadeau-Bengio corrected $p = \CorrectedTP$). This parity confirms that 2D connectivity matrices capture the requisite pharmacophoric variance.

Parsimony therefore drives the optimal screening strategy. Excluding 3D descriptors and synthetic data augmentation eliminates processing latency during both feature extraction and model training. For this specific target, high-dimensional representations introduce computational friction without improving model fidelity. Consequently, lightweight 2D fingerprint ensembles offer a strictly superior, highly scalable architecture for virtual screening under constrained data conditions.

% ==========================================
% 4. SECTION 4.3
% ==========================================

\subsection{Methodological Scope, Limitations, and Translational Potential}

The predictive reliability of any quantitative structure-activity relationship (QSAR) model is inherently bounded by the chemical diversity of its training corpus. Given the dataset size of \NCompounds\ compounds, the applicability domain of our optimized ensemble is strictly restricted to the structural scaffolds and physicochemical properties represented within the curated TgDHFR data. Consequently, extrapolating predictions to highly divergent chemotypes or novel synthetic spaces carries an increased risk of reduced predictive confidence. 

A specific vulnerability intrinsic to 2D circular fingerprints in low-data regimes is their susceptibility to activity cliffs—pairs of structurally similar molecules that exhibit radically divergent bioactivities due to subtle chemical modifications. A single bioisosteric substitution or methylation can alter a critical hydrogen bond or steric fit within the binding pocket. Because Morgan fingerprints map topological neighborhoods into bit vectors without explicitly encoding 3D spatial conformation or electrostatic field potentials, they risk mapping active and inactive cliff-forming analogs to virtually identical representations. We strategically mitigated this limitation through our downstream multi-tier architecture. The orthogonal integration of physics-based molecular docking explicitly evaluates 3D spatial complementarity and steric clashes, acting as a structural filter that detects activity cliffs otherwise invisible to pure 2D topological models.

An additional methodological constraint involves the extensive architecture search executed within the inner optimization loop. Evaluating up to 41 distinct stacking combinations across six regressor families introduces a latent risk of model selection bias. While the nested cross-validation framework mathematically isolates this optimization phase from the outer test folds—thereby mitigating spurious performance inflation—it does not entirely eliminate architectural instability. To transparently quantify this variance, Supplementary Table S1 details the selection frequency of each meta-ensemble across the \TotalOuterFolds\ outer folds. Within the optimal Morgan fingerprint space, the \MorganBestArchitecture\ configuration emerged as the most stable architecture, yet it was selected in only \MorganStability\% of the iterations. This distribution confirms that while the nested design successfully prevents test-set leakage, the inherent data sparsity of the \NCompounds-compound benchmark inevitably restricts the convergence of the inner loop onto a single, universally dominant stacking hierarchy.

Despite these boundary conditions, the programmatic pipeline developed in this study exhibits high transferability to other neglected tropical diseases (NTDs) and orphan targets. Because the framework relies on classical ensemble learning rather than data-hungry deep neural networks, it maintains predictive stability in low-data regimes (e.g., datasets with fewer than 1,000 compounds). This architecture provides a computationally lightweight triage mechanism that can be rapidly deployed for other under-researched pathogens. By systematically filtering extensive libraries of approved drugs to isolate a tractable subset of candidates, the pipeline minimizes the computational overhead required for subsequent physics-based simulations and reduces the attrition rate in downstream experimental validations.

Finally, the integration of ADMET profiling contextualizes the clinical viability of the prioritized candidates. The top-ranked hybrid hit, Pralatrexate (\DockBestDockScore\ kcal/mol), alongside the reference drug Methotrexate, exhibited low predicted Caco-2 permeability, classifying them as parenteral rather than oral candidates. Rather than a limitation, this pharmacokinetic profile is mechanistically consistent with classical antifolates, which are typically highly polar, structurally complex molecules. Clinically, parenteral administration aligns directly with the treatment protocols for severe, acute manifestations of the disease, such as toxoplasmic encephalitis in immunocompromised patients, where intravenous delivery ensures rapid systemic bioavailability. Consequently, the pipeline not only identifies potent target binders but also accurately anticipates their clinical administration routes, reinforcing its utility in applied pharmacological repurposing.

